\documentclass[10pt,a4paper]{amsart}
\usepackage[margin=19mm]{geometry}
\usepackage{amsmath,amssymb,mathtools}
\usepackage[scr=rsfs,cal=euler]{mathalfa}
\usepackage{xfrac}
\usepackage[unicode,hidelinks,bookmarks=false]{hyperref}
\newcommand{\ie}{i.e.\ }
\newcommand{\cf}{cf.\ }
\newcommand{\resp}{respectively\ }
\newcommand{\Subsection}{\subsection}
\providecommand{\nobreakdash}{\nobreak\hbox{-}}
\setlength{\emergencystretch}{2em}
\multlinegap0pt
\usepackage{amssymb}
\newtheorem{df}{Definition}
\newtheorem{cor}{Corollary}
\newtheorem{thm}{Theorem}
\newtheorem{lm}{Lemma}
\newcommand{\cn}[1]{{\con\m {#1}}}
\newcommand{\con}{{\rm Con\:}}
\newcommand{\lb}{\langle}
\newcommand{\m}[1]{{\uppercase {\mathbf {#1}}}}
\newcommand{\mt}{\wedge}
\newcommand{\rb}{\rangle}
\newcommand{\usub}{\subseteq}
\newcommand{\vr}[1]{{\uppercase {\mathcal{#1}}}}
\title{SMB algebras II: On the Constraint Satisfaction Problem over Semilattices of Mal'cev Blocks}
\author{Petar Markovi\'c, Mikl\'os Mar\'oti, Ralph McKenzie, Aleksandar Proki\'c}
\date{Source: arXiv:2604.05161v1 (CC BY 4.0)}
\begin{document}
\maketitle

\noindent\emph{Source and licence.} SMB algebras II: On the Constraint Satisfaction Problem over Semilattices of Mal'cev Blocks. Version arXiv:2604.05161v1 (CC BY 4.0), distributed by arXiv under the Creative Commons Attribution 4.0 International licence, http://creativecommons.org/licenses/by/4.0/, as recorded in the arXiv licence metadata for this exact version and shown on the abstract page https://arxiv.org/abs/2604.05161v1. Full licence notice: \texttt{licenses/arxiv-2604.05161-CC-BY-4.0.txt}.\par\smallskip

\noindent\textit{Editorial scope.} Counted unit: the article's thm flattract (source0.tex lines 309--311). The article's complete original proof of it is delivered with it (source0.tex lines 313--401, one \texttt{proof} environment of 89 lines). No mathematics is written, repaired or re-derived: the statement and the proof are the article's own lines, in the article's own notation, and no retained line is substituted or re-flowed.  Retained uncounted as the source-local context the proof needs: lines 229--237; lines 243--245.  Nothing was omitted from any retained span, so the package carries no omission record.  No retained line contains a citation, so the wrapper carries no bibliography and no bibliography entry.  No retained line contains a figure environment, an image-inclusion command or drawing code, so no image is delivered and the package declares no asset dependency.  Editorial formatting only, in addition: an AMS \texttt{amsart} preamble that restates the article's own declarations and macros used by the retained lines, namely the environments \texttt{df}, \texttt{cor}, \texttt{thm}, \texttt{lm} on separate counters, together with the standard packages those lines need.  The retained environments print in this document as Df 1, Corollary 1, Theorem 1, Lm 1 in order of appearance, whereas the article prints the same items with the numbering of its own full document.  The complete native source of the article as distributed in the arXiv e-print for this exact version is bundled unchanged as \texttt{sources/197933/source0.tex}.
\begin{df}\label{regSMBdef}
We say that an SMB algebra $\m a=\lb A;\mt,d\rb$ over ${\sim}\in\cn a$ is {\em regular} if 
\begin{enumerate}
\item\label{property1} for all $a,b,c\in A$, $[d(a,b,c)]_{\sim}=[(a\mt b)\mt c]_{\sim}$,
\item for all $a,b\in A$ such that $[b]_{\sim}\geq[a]_{\sim}$, $a\mt b=a$,
\item $\m a\models d(x,y,z)\approx d(x\mt (z\mt y),y\mt (z\mt x),z\mt (y\mt x))$, and
\item $\m a\models x\mt (x\mt y)\approx x\mt y$.
\end{enumerate}
\end{df}

\begin{cor}\label{regSMBsubalg}
If $S$ is any subuniverse of the semilattice $\m a/{\sim}$ induced by the regular SMB algebra $\m a$, then the union of all $\sim$-classes in $S$ is a subuniverse of $\m a$. In particular, the conclusion holds if $S$ is a down-set (order ideal) or an interval in the partially ordered set $(A/{\sim};\leq)$.
\end{cor}

\begin{thm}\label{flattract}
Let $\m a=\lb A;w\rb$ be an SMB algebra wrt. $\sim$. $CSP(\m a)$ is tractable when the order of $\sim$-classes is that of a flat semilattice.
\end{thm}

\begin{proof}
We assume, as before, that $\m a$ is a regular SMB algebra and ensure that $P$ is $(2,3)$-minimal so we may call $P=([n],D,\vr c)$ a multisorted instance of $CSP(\vr t(\m a))$ whose domains are subuniverses of $\m a$. Let $A=O\cup I_1,\ldots,I_m$, where $O$ is the least $\sim$-class and $I_1,\ldots,I_m$ are the maximal ones. 

Because of $(2,3)$-minimality, not only the domains $A_1$ are fixed by the instance, but also $P$ contains all constraints of the form $(\{i,i'\},R_{i,i'})$, and $R_{i,i'}$ is the projection to $\{i,i'\}$ of any constraint relation $R$ such that $(S,R)\in \vr c$ and $\{i,i'\}\subseteq S$. We think of $R_{i,i'}$ as a bipartite graph between disjoint partitions $A_i$ and $A_{i'}$, and the minimality conditions imply that there are no isolated vertices.

$P$ defines a preorder $\preceq$ on the following set: 
$$S(P):=\{(i,j):0\leq i<n \;\&\; 1\leq j\leq m \;\&\; A_i\cap I_j\neq\emptyset\neq O\cap A_i\}.$$ 
We say $(i,j)\preceq(i',j')$ if, whenever $f\in R_{i,i'}$ and $f(i)\in I_j$, then $f(i')\in I_{j'}$ (here $(\{i,i'\},R_{i,i'})\in\vr c$). We define $(i,j)\cong(i',j')$ if $(i,j)\preceq(i',j')$ and $(i',j')\preceq(i,j)$. We call $\cong$-classes {\em strands}.

It is not hard to see that $\preceq$ is indeed a preorder: reflexivity is a trivial consequence of the definition of $\preceq$. For transitivity, assume that $(i,j)\preceq(i',j')$ and $(i',j')\preceq(i'',j'')$, and let $f\in R_{i,i''}$ be such that $f(i)\in I_j$. By $(2,3)$-minimality, there exists $\overline{f}\in R_{i,i',i''}$ such that $\overline{f}|_{\{i,i''\}}=f$. Denote $g=\overline{f}|_{\{i,i'\}}$ and $h=\overline{f}|_{\{i',i''\}}$. From $g(i)=\overline{f}(i)=f(i)\in I_j$ and $(i,j)\preceq(i',j')$ follows $h(i')=\overline{f}(i')=g(i')\in I_{j'}$, similarly, from $(i',j')\preceq(i'',j'')$ we obtain $f(i'')=h(i'')\in I_{j''}$.

The properties to remember about strands are (all of these follow from $(2,3)$-minimality of $P$): Firstly, for any strand and $i\leq n$, there is at most one $(i,j)$ in this strand (the argument uses that the meet of any two elements of $A_i$ which are in different $\sim$-classes must be in $O$). Secondly, for any pair $(i,j)\in S(P)$, it is in exactly one strand since $\cong$ is an equivalence relation. Finally, whenever $f(i)\in I_j$ for some solution $f\in A^n$ of $P$ (or of a restriction of $P$ to some subset of variables), then for all $i'\in[n]$ such that $(i',j')$ is in the same strand as $(i,j)$, then $f(i')\in I_{j'}$ (in case of a restriction of $P$, we need to assume here that $i'$ is in the scope of this restriction, as well). So, we may refer to a strand as $\{(i,t(i)):i\in F\}$ where $F\usub [n]$ and $t:F\rightarrow \{1,\ldots,m\}$.

Now we describe a procedure $TESTASTRAND(P,\{(i,t(i)):i\in F\})$.

\begin{enumerate}
\item[{\bf Step 1.}] Compute the restriction $Q:=P|_F$.
\item[{\bf Step 2.}] Using the Mal'cev algorithm decide if $Q$ has an $O$-valued solution.
\begin{enumerate}
\item[{\bf Step 2.1.}] If 'YES' output 'Exists an $O$-valued solution' and stop.
\item[{\bf Step 2.2.}] If 'NO', test if $Q$ has a solution $f$ which for all $i\in F$ has $f(i)\in t(i)$. This is again a Mal'cev algorithm. If 'YES', output 'Exists a strand-valued solution' and stop, if 'NO', output 'Neither' and stop.
\end{enumerate}
\end{enumerate}

We will repeatedly use this procedure. The main algorithm $SOLVE(P)$ now looks like this:

\begin{enumerate}
\item[{\bf Step 1.}] Set up the working instance $Q:=P$.
\item[{\bf Step 2.}] Applying consistency algorithm reduce $Q$ to a $(2,3)$-minimal instance. If any $(S,R)\in\vr c(Q)$ is such that $R=\emptyset$, output '$P$ has no solution' and stop.
\item[{\bf Step 3.}] Compute $S(Q)$ and $J=\{i\in n:(\exists j)(i,j)\in S(Q)\}$. Compute $\preceq$ and $\cong$ on $S(Q)$.
\item[{\bf Step 4.}] For each strand $\{(i,t(i)):i\in F\}$ do
\begin{enumerate}
\item[{\bf Step 4.1.}] If $TESTASTRAND(Q,\{(i,t(i)):i\in F\})$ outputs 'Exists an $O$-valued solution', move on to next strand.
\item[{\bf Step 4.2.}] If  $TESTASTRAND(Q,\{(i,t(i)):i\in F\})$ outputs 'Exists a strand-valued solution', tighten $Q$ by replacing each $A_i$, $i\in F$, with $I_{t(i)}$ and go to Step 2.
\item[{\bf Step 4.3.}] If $TESTASTRAND(Q,\{(i,t(i)):i\in F\})$ outputs 'Neither', tighten $Q$ by replacing each $A_i$, $i\in F$, with  $A_i\setminus I_{t(i)}$ and go to Step 2.
\end{enumerate}
\item[{\bf Step 5.}] Tighten $Q$ by replacing each $A_i$, $i\in J$, with  $A_i\cap O$.
\item[{\bf Step 6.}] $Q$ is now an instance over Mal'cev algebras, as each $A_i$ is now within a single $\sim$-class, so solve it. Output that $P$ has a solution if $Q$ does and that it has no solution if $Q$ does not and stop.
\end{enumerate}

We see that the instance always tightens to a smaller one whenever either Step 4.2 or Step 4.3 is applied. Therefore, they can be applied at most $n|A|$ many times. Moreover, the number of strands is bounded from above by $nm\leq n|A|$, so at most $n^2|A|^2$ many applications of the TestAStrand procedure can occur during running time of the algorithm. Each application requires at most two applications of the Mal'cev algorithm. Moreover, we need to apply the consistency algorithm in Step 2 at most $n|A|+1$ times. Finally, there is one more application of the Mal'cev algorithm in Step 6. So the time-cost of the algorithm is dominated by the number of steps needed to apply $(2,3)$-minimality algorithm $n|A|+1$ times, plus the number of steps needed to apply the Mal'cev algorithm $2n^2|A|^2+1$ times. As both of these elementary algorithms run in  polynomial time, so does this one.

To study the correctness of our algorithm, we note first that when it outputs that there exists a solution to $P$, then the algorithm really found a solution to a tightening of $P$ to certain specified path through $\sim$-classes (since when the Mal'cev algorithm produces 'YES' it actually finds a generating set for the space of all such solutions). Therefore, this is a solution to $P$, as well.

If the algorithm rejected the instance, on the other hand, it may have done so in the Step 2., or in Step 6. If it rejected $P$ in Step 2., it either did it in the initial $2$-consistency check, which is fine as the instance which fails the initial consistency check indeed has no solution, or in application of Step 2. which occurred after a tightening via either Step 4.2. or 4.3.

Both of these tightenings lead to equivalent instances. To see this, fix a strand $\{(i,t(i)):i\in F\}$. If the instance $P|_F$ has a solution which at some (equivalently, all) $i\in F$ takes values in $I_{t(i)}$ and another one which at some (equivalently, all) $i\in F$ takes values outside the $I_{t(i)}$, then the meet of these two solutions is a solution of $P|_F$ which is $O$-valued at each variable $i\in F$ (recall that by definition of $S(P)$, for all $i\in F$, $A_i\cap O\neq\emptyset$). Therefore, if there are no $O$-valued solutions to $P|_F$, then either $P|_F$ has no strand-valued solutions, or it has no solutions which are outside the strand at all $i\in F$. Again, by the definition of a strand, these are the only two options. Note that the relation $\cong$ is invariant under restriction since $P$ is $(2,3)$-minimal, so $\{(i,t(i)):i\in F\}$ is a strand for $P|_F$, as well as for $P$. Also note, by Corollary~\ref{regSMBsubalg}, that both $I_{t(i)}$ and $A_i\setminus I_{t(i)}$ are subuniverses of $\m a$, so the tightened instances are still within $CSP(\vr t(\m a))$.

Finally, the algorithm may have rejected the instance at Step 6. To see that this is a correct rejection, we need to prove the following

\begin{lm}\label{Malcev}
Assume that $f$ is a solution to a $(2,3)$-minimal multisorted instance $P=([n],D,\vr c)$ of $CSP(\vr t(\m a))$. If for all strands $\{(i,t(i)):i\in F\}$ such that $f(i)\in I_{t(i)}$ there exists an $O$-valued solution of $P|_F$, then $P$ has a solution $\overline{f}\in A^n$ such that $\overline{f}\in O$ whenever $A_i\cap O\neq\emptyset$.
\end{lm}

\begin{proof}
We first note that, since $f$ is a function, for any two different strands $\{(i,t_1(i)):i\in F_1\}$ and $\{(i,t_2(i)):i\in F_2\}$ visited by $f$ (i. e. such that $f(i)\in I_{t_1(i)}$ when $i\in F_1$ and $f(i)\in I_{t_2(i)}$ when $i\in F_2$), we must have that $F_1$ and $F_2$ are disjoint. Therefore, we know that the set $J\usub n$ defined by $J=\{i\in n:f(i)\notin O\;\&\;A_i\cap O\neq\emptyset\}$ is partitioned into $F_1\cup\ldots\cup F_l$, which are the sets of variables of all strands visited by $f$. We may define now uniformly $t:J\rightarrow\{1,\ldots,m\}$ such that for $i\in F_j$, $t(i):=t_j(i)$. As the strands are partially ordered by $\preceq$, we assume without loss of generality that when $i\in F_j$, $i'\in F_{j'}$ and $1\leq j'<j\leq l$, then $\neg(i,t_j(i))\preceq(i',t_{j'}(i'))$. Our assumptions give us $O$-valued solutions $h_j$ to $P|_{F_j}$, for all $1\leq j\leq l$. We prove inductively on $s$ that there exists a solution $g_s$ of $Q_s=P|_{(n\setminus J)\cup F_1\cup\ldots\cup F_s}$ which is equal to $f$ on $n\setminus J$ and to $f\mt h_j$ (which is $O$-valued) on each $F_j$, $1\leq j\leq s$. Clearly for $s=l$ this will be the desired $\overline{f}$, as $Q_l=P$.

The base case is for $s=1$, where we claim that $g_1(i)=f(i)$ for $i\in n\setminus J$ and $g_1(i)=f(i)\mt h_1(i)$ is a solution to $Q_1$. Indeed, let $(S,R)\in\vr c(Q_1)$. If $S$ is disjoint from $F_1$, then $g_1|_S=f|_S$, so $g_1|_S\in R$. If, on the other hand, $S$ is not disjoint from $F_1$, then from the fact that $h_1$ is a solution to $P|_{F_1}$, we know that there exists some $h\in R$ such that $h|_{S\cap F_1}=h_1|_{S\cap F_1}$. Therefore, where $f'=f|_S$, $f'\mt h\in R$. Now for $i\in S\setminus F_1$, $f'(i)\in\mathrm{min}(\m A_i)$, so $f'(i)\mt h(i)=f'(i)=g_1(i)$. On the other hand, where $i\in S\cap F_1$, $f'(i)\mt h(i)=f(i)\mt h_1(i)=g_1(i)$. Therefore, $g_1$ is a solution to $Q_1$.

Assume that $g_s$ is a solution to $Q_s$, where $g_s$ is
defined as above. Let $(S,R)\in \vr c(Q_{s+1})$. If $S\cap F_{s+1}=\emptyset$, there is nothing to prove, and if $S\cap (F_1\cup\ldots\cup F_s)=\emptyset$, this is the same as the base case $s=1$. So we assume that $S$ intersects both $F_1\cup\ldots\cup F_s$ and $F_{s+1}$.

We need to construct some auxiliary elements of $R$. Firstly, we are guaranteed to have $f',g',h'\in R$ such that $f'=f|_S$, $g'|_{S\setminus F_{s+1}}=g_s|_{S\setminus F_{s+1}}$, while $h'|_{S\cap F_{s+1}}=h_{s+1}|_{S\cap F_{s+1}}$, where the existence of $g'$ and $h'$ follows from the definition of the restriction of an instance to a set of coordinates. Secondly, we define $S_0=S\setminus(F_1\cup\ldots\cup F_{s+1})$ and $S_j=S\cap F_j$ for $1\leq j\leq s+1$. Our assumptions ensure that $S_{s+1}\neq\emptyset$ and for at least one $j$, $1\leq j\leq s$, $S_j\neq\emptyset$. Note also that $S_0$ consists precisely of those variables $i\in S$ where $f(i)\in\mathrm{min}(A_i)$, so $f(i)$ is a left absorbing element for $\mt$ on $A_i$ (here we use the fact that $(S,R)\in \vr c(Q_{s+1})$, so $S$ has no variables in $F_j$, $j>s+1$).

Now select some $i_{s+1}\in S_{s+1}$ and for any $j$ such that $1\leq j\leq s$ and $S_j\neq \emptyset$ a variable $i_j\in S_j$. Let $(\{i_j,i_{s+1}\},R_j)\in \vr c(Q_{s+1})$. As we know that $\neg(i_{s+1},t(i_{s+1}))\preceq(i_j,t(i_j))$ and that $(f(i_j),f(i))\in R_j\cap (I_{t(i_j)}\times I_{t(i_{s+1})})$, then there must exist also some $q_j\in R_j$ such that $q_j(i_j)\notin I_{t(i_j)}$, while $q_j(i_{s+1})\in I_{t(i_{s+1})}$. Let $q_j'\in R$ be such that $q_j'|_{\{i_j,i_{s+1}\}}=q_j$ and $p_j'=f'\mt q_j'$. We know that $q_j'|_{S_j}(i)\notin I_{t(i)}$ for all $i\in S_j$, as it is so at $i_j$ and $\{(i,t(i)):i\in S_j\}$ is the restriction of a strand. Therefore, for all $i\in S_j$, $p_j'(i)\in O$. Moreover, for all $i\in S_0\cup S_{s+1}$, $p_j'(i)=f(i)$. Now define $p'=\bigwedge\limits_{S_j\neq \emptyset}p_j'$ to be the left-associated meet. We know that $p'\in R$, that $p'|_{S_0\cup S_{s+1}}=f|_{S_0\cup S_{s+1}}$ and that $p'(i)\in O$ for all $i\in S_j$, $1\leq j\leq s$.

Now we define three tuples $f_1$, $f_2$ and $f_3$ in $R$ in the following way: $$f_1=f'\mt (g'\mt h')\text{, }f_2=p'\mt f_1\text{ and }f_3=p'\mt h'.$$ We claim that $\overline{f}|_S=d(f_1,f_2,f_3)$, thus proving $\overline{f}|_S\in R$. We break down the proof into three cases: 

If $i\in S_0$, then $\overline{f}(i)=f(i)=f'(i)=p'(i)\in\mathrm{min}(\m A_i)$, so it is a left absorbing for the meet. Therefore, $f_1(i)=f_2(i)=f_3(i)=\overline{f}(i)$, and the claim follows by the idempotence of $d$.

If $i\in S_j$ for some $1\leq j\leq s$, then $f_2(i)=p'(i)=f_3(i)\in O$, while 
$$
\begin{gathered}
f_1(i)=f(i)\mt (g'(i)\mt h'(i))=f(i)\mt g'(i)=\\
f(i)\mt(f(i)\mt h_j(i)) =f(i)\mt h_j(i)=\overline{f}(i)\in O
\end{gathered}
$$ (the second equality follows from $g'(i)\in O$, while the fourth equality follows from Definition~\ref{regSMBdef} (4)). Therefore, $$d(f_1,f_2,f_3)(i)=d(f_1(i),p'(i),p'(i))=f_1(i)=\overline{f}(i),$$ since $d$ is a Mal'cev operation on $O$.

Finally, if $i\in S_{s+1}$, then $g'(i)\mt h_{s+1}(i)\in O$ and $$f_1(i)=f(i)\mt(g'(i)\mt h_{s+1}(i)),$$ while $p'(i)=f(i)$ and therefore 
$$
\begin{gathered}
f_2(i)=f(i)\mt f_1(i)=f(i)\mt (f(i)\mt(g'(i)\mt h_{s+1}(i)))=\\
f(i)\mt(g'(i)\mt h_{s+1}(i))=f_1(i).
\end{gathered}
$$ 
On the other hand, $$f_3(i)=p'(i)\mt h_{s+1}(i)=f(i)\mt h_{s+1}(i)=\overline{f}(i)\in O.$$ Now $$d(f_1,f_2,f_3)(i)=f_3(i)=\overline{f}(i),$$ since $d$ is a Mal'cev operation on $O$. This finishes the inductive proof of the lemma.
\end{proof}
We have proved that, if there exists a solution of an instance and if for each strand there exists an $O$-valued solution to the restriction of the instance to the strand variables, then there exists an $O$-valued solution of the whole instance. If the instance $Q$ reached Step 6, there is an $O$-valued solution to the restriction of $Q$ to the strand variables, for each strand, hence testing $Q$ just for $O$-valued solutions determines whether $Q$ has a solution. 
\end{proof}

\end{document}
